To the Members of the WHO Scientific Advisory Group for the Origins of Novel Pathogens (SAGO)  
and to Drs. Havens, Kosakovsky Pond, Zehr, Pekar, Parker, Worobey, Andersen, Wertheim and colleagues:

I write to draw your attention to a publicly available technical analysis that formalizes a chronological constraint relevant to the proximal-origin question for SARS-CoV-2.

The materials are deposited at:  
https://github.com/CovidResearcher/COVID-THE-CODE

The central documents are:

the_6_month_window.tex  
Covid_Timeline.py  
proximal_origin.tex  
tenable_threshold.tex  
side_by_side_comparative_chart.tex  
(and the supporting files hyper_mutation_fcs_aquisition_steps.tex, june_december_2019_zoonosis_deadline.tex)

These files set out the following argument in quantitative and chronological terms:

Documented intensive sampling of the relevant southern Chinese and Southeast Asian bat reservoirs effectively closed on 31 May 2019.  
From that date until the first detectable human cases in Wuhan lies a window of approximately six months (June–December 2019).  
Within that interval a pure natural-proximal account must either recover a specimen carrying the exact 12-nucleotide PRRA insertion with the rare CGG-CGG codon pair, or explain its absence while still generating a fully human-adapted virus that appears, fully formed, at a single geographic point.  
No such natural intermediate has been recovered. Subsequent sampling continues to yield nearer relatives that still lack the precise configuration.  
Molecular-clock considerations render later material incompatible with a temporally authentic 2019 proximal ancestor; evolutionary time already elapsed is irreversible.  
The alternative—positing a rapid, multi-host cascade of recombination, fixation of the exact furin-cleavage insertion, and silent amplification inside six months—requires an extraordinary sequence of low-probability events that left no recoverable intermediates.

A side-by-side comparative chart in the repository contrasts the requirements of that multi-species natural cascade against a research-related continuum that does not demand the same unobserved intermediate steps inside the same narrow window.

I submit these materials for your consideration as part of the ongoing evaluation of hypotheses concerning the origin of SARS-CoV-2. The chronological compression after the close of intensive reported sampling is, in my view, decisive and has not been adequately addressed in the published literature to date.

The full repository, including the executable timeline scripts and the formal LaTeX expositions, is ready for mathematical and immunovirological confirmation.

To the Members of the WHO Scientific Advisory Group for the Origins of Novel Pathogens (SAGO)  
and to Drs. Havens, Kosakovsky Pond, Zehr, Pekar, Parker, Worobey, Andersen, Wertheim and colleagues:

I write to draw your attention to a publicly available technical analysis that formalizes a chronological and molecular-clock constraint relevant to the proximal-origin question for SARS-CoV-2.

The materials are deposited at:  
https://github.com/CovidResearcher/COVID-THE-CODE

thesis.tex states the core claim directly:

“The practical possibility of recovering a temporally matching 2019 specimen from extant colonies has already reached zero.  
This is not a rhetorical flourish or an administrative cutoff. It is a direct, irreversible consequence of molecular evolution operating under ordinary biological conditions.”

The argument rests on three linked points:

Continuous action of the molecular clock  
   RNA viruses accumulate substitutions at measurable rates (typically \(10^{-3}\) to \(10^{-4}\) substitutions per site per year in structural genes). Any lineage that existed in 2019 and continued circulating through 2020–2026 has necessarily experienced six to seven additional years of mutational accumulation. A specimen recovered in 2026 therefore belongs to a descendant generation and cannot be genetically identical—or sufficiently proximal—to the precise genome required as the immediate pre-spillover ancestor of late 2019.

Absence of the defining configuration  
   No natural specimen carrying the exact polybasic furin cleavage site (PRRA insertion with the CGG-CGG codon pair) has been documented from the intensive sampling that ran through 31 May 2019 or from subsequent years. The nearest relatives lack the motif entirely.

Temporal irreversibility  
   Living reservoirs supply only contemporary material. They cannot furnish a frozen 2019 genome. Future sampling may recover closer overall relatives or independent furin-like motifs, but it cannot recover a virus that stopped evolving in 2019 and waited for discovery. The window in which a temporally matching natural proximal ancestor could have been recovered closed with the passage of calendar time itself.

thesis.tex concludes that pure natural-proximal accounts must therefore abandon hope of recovering a temporally authentic specimen and are left with residual options that multiply unobserved entities (entirely unreported local populations that later went extinct or evolved beyond recognition, or an immediate low-probability multi-species cascade that left no recoverable intermediates). These stand in contrast to a research-related continuum that does not require a hidden natural proximal ancestor circulating outside the documented field-to-lab pathway.

I submit these materials, and especially the formal exposition in thesis.tex, for your consideration. The evolutionary-clock barrier is presented not as proof of any single origin hypothesis, but as a physical constraint that permanently tightens the requirements on natural-proximal narratives.

The full repository is open for inspection and critique.

Respectfully,  
Covid Researcher

These three points form a single, interdependent argument about the practical impossibility of recovering a temporally authentic natural proximal ancestor of SARS-CoV-2 from living animal populations after several years of continued evolution.

1. Continuous action of the molecular clock

RNA viruses replicate with an error-prone RNA-dependent RNA polymerase. This generates a measurable background rate of nucleotide substitutions, insertions, deletions, and recombination events while the virus circulates in any living host.

Empirical estimates for coronaviruses, including sarbecoviruses, place the substitution rate in structural genes on the order of \(10^{-3}\) to \(10^{-4}\) substitutions per site per year. These rates are not arbitrary; they are the same calibrated clocks routinely applied in Bayesian phylogenetic analyses (BEAST and related frameworks) to date the divergence of RaTG13, the BANAL viruses, RmYN02, and other published SARS-related sequences against known sampling dates.

Because the clock operates continuously whenever the virus is replicating, any lineage that existed in late 2019 and continued circulating in bats, intermediate hosts, or other reservoirs through 2020, 2021, 2022, 2023, 2024, 2025, and into 2026 necessarily accumulated six to seven additional years of mutational change.  

A specimen collected from those populations in 2026 therefore belongs to a descendant generation. It differs from its 2019 progenitor by dozens to hundreds of substitutions (depending on exact rate, generation time, and population size), plus possible recombination events. It cannot be genetically identical—or even sufficiently proximal in the strict phylogenetic sense required for a direct pre-spillover ancestor—to the precise genome that would have been circulating immediately before the late-2019 emergence of SARS-CoV-2.

Evolution does not pause, reverse, or hold a lineage in stasis simply because later investigators desire a pristine temporal match. The molecular clock is a physical process, not an administrative one.

2. Absence of the defining configuration

The defining molecular feature under discussion is the precise polybasic furin cleavage site: the 12-nucleotide insertion that encodes the amino-acid motif PRRA, specifically with the rare tandem arginine codons CGG-CGG, located at the S1/S2 junction of the spike protein.

No natural specimen carrying this exact configuration has been documented from the intensive sampling campaigns that ran through 31 May 2019, nor from the subsequent years of expanded surveillance in southern China, Southeast Asia, or other regions.  

This absence is not merely a gap in sampling effort; it is the empirical starting condition against which the molecular-clock argument is applied. If the exact configuration had been present in surveyed reservoirs before or during mid-2019, the documented scale of the sampling network (thousands of bats, multiple provinces, full-genome and partial-gene sequencing) makes its non-recovery increasingly difficult to attribute solely to incomplete coverage.

3. Temporal irreversibility

Living animal reservoirs can supply only contemporary viral genomes. They cannot supply a “frozen” 2019 genome.

Suppose, for the sake of argument, that an undetected natural lineage carrying the exact 2019 configuration (PRRA + CGG-CGG) had been circulating in some unsampled colony or trade chain in 2019. That lineage would have continued to evolve under the same molecular clock. By 2026 its descendants would already differ from the 2019 progenitor by the cumulative mutational distance described above.  

Recovering a modern specimen and claiming it represents the missing 2019 proximal ancestor would therefore require either:
that the lineage somehow ceased evolving for six-plus years (biologically impossible under known RNA-virus dynamics), or  
that investigators could rewind the molecular clock (physically impossible).

Neither option exists. Future expeditions into the same karst systems, the same Rhinolophus colonies, or expanded wildlife-trade networks may recover closer overall relatives, additional recombinants, or even independent furin-like motifs in distant lineages. What they cannot recover is a virus that stopped evolving in 2019 and waited patiently for discovery.

The window in which a temporally matching natural proximal ancestor could have been recovered therefore closed with the simple passage of calendar time. Intensified future sampling cannot reopen that window; each additional year of negative results further reduces the probability that the exact form was simply overlooked, while the evolutionary distance between any living population and the 2019 emergence window continues to increase.

Combined implication

Taken together, the continuous molecular clock, the documented absence of the exact PRRA + CGG-CGG configuration, and the irreversibility of elapsed evolutionary time establish a physical constraint: the practical possibility of recovering a temporally authentic 2019 natural proximal ancestor from extant colonies has already reached zero.  

This does not, by itself, prove any particular origin hypothesis. It does permanently tighten the requirements on pure natural-proximal accounts, which must now rely on residual scenarios that multiply unobserved entities (completely unreported local populations that later went extinct or evolved beyond recognition, or an immediate low-probability multi-species cascade that left no recoverable intermediates).  

The argument is presented as a consequence of ordinary viral population genetics and the arithmetic of elapsed time, not as a temporary gap that further fieldwork can close.
